\section*{NeurIPS Paper Checklist}

%%% BEGIN INSTRUCTIONS %%%
The checklist is designed to encourage best practices for responsible machine learning research, addressing issues of reproducibility, transparency, research ethics, and societal impact. Do not remove the checklist: {\bf The papers not including the checklist will be desk rejected.} The checklist should follow the references and follow the (optional) supplemental material.  The checklist does NOT count towards the page
limit. 

Please read the checklist guidelines carefully for information on how to answer these questions. For each question in the checklist:
\begin{itemize}
    \item You should answer \answerYes{}, \answerNo{}, or \answerNA{}.
    \item \answerNA{} means either that the question is Not Applicable for that particular paper or the relevant information is Not Available.
    \item Please provide a short (1--2 sentence) justification right after your answer (even for \answerNA). 
   % \item {\bf The papers not including the checklist will be desk rejected.}
\end{itemize}

{\bf The checklist answers are an integral part of your paper submission.} They are visible to the reviewers, area chairs, senior area chairs, and ethics reviewers. You will also be asked to include it (after eventual revisions) with the final version of your paper, and its final version will be published with the paper.

The reviewers of your paper will be asked to use the checklist as one of the factors in their evaluation. While \answerYes{} is generally preferable to \answerNo{}, it is perfectly acceptable to answer \answerNo{} provided a proper justification is given (e.g., error bars are not reported because it would be too computationally expensive'' or ``we were unable to find the license for the dataset we used''). In general, answering \answerNo{} or \answerNA{} is not grounds for rejection. While the questions are phrased in a binary way, we acknowledge that the true answer is often more nuanced, so please just use your best judgment and write a justification to elaborate. All supporting evidence can appear either in the main paper or the supplemental material, provided in appendix. If you answer \answerYes{} to a question, in the justification please point to the section(s) where related material for the question can be found.

IMPORTANT, please:
\begin{itemize}
    \item {\bf Delete this instruction block, but keep the section heading ``NeurIPS Paper Checklist"},
    \item  {\bf Keep the checklist subsection headings, questions/answers and guidelines below.}
    \item {\bf Do not modify the questions and only use the provided macros for your answers}.
\end{itemize} 
 

%%% END INSTRUCTIONS %%%


\begin{enumerate}

\item {\bf Claims}
    \item[] Question: Do the main claims made in the abstract and introduction accurately reflect the paper's contributions and scope?
    \item[] Answer: \answerYes{}
    \item[] Justification: The abstract and \S\ref{sec:intro} state three claims: (i) DISCA is a train-free inference-time controller that requires no per-country reward models or extra training data; (ii) it cuts macro MIS by 10--24\% on seven open-weight backbones across 20 countries and 6 moral dimensions; (iii) calibration competes with scale (Phi-4 14B beats Llama-3.3-70B in absolute MIS). All three are substantiated in Table~\ref{tab:main_macro_summary} and Figure~\ref{fig:scaling_vs_calibration}; scope and assumptions (binary dilemmas with decision-token logits, WVS-7 vintage) are made explicit in \S\ref{sec:discussion} and Appendix~\ref{app:limitations}.
    \item[] Guidelines:
    \begin{itemize}
        \item The answer \answerNA{} means that the abstract and introduction do not include the claims made in the paper.
        \item The abstract and/or introduction should clearly state the claims made, including the contributions made in the paper and important assumptions and limitations. A \answerNo{} or \answerNA{} answer to this question will not be perceived well by the reviewers. 
        \item The claims made should match theoretical and experimental results, and reflect how much the results can be expected to generalize to other settings. 
        \item It is fine to include aspirational goals as motivation as long as it is clear that these goals are not attained by the paper. 
    \end{itemize}

\item {\bf Limitations}
    \item[] Question: Does the paper discuss the limitations of the work performed by the authors?
    \item[] Answer: \answerYes{}
    \item[] Justification: \S\ref{sec:discussion} bounds generalisation with three caveats (alignment to a survey statistic vs.\ perceived legitimacy; PT value function as aggregation kernel rather than cognitive model; decision-token logit requirement). Appendix~\ref{app:limitations} expands these with: WVS-to-trolley linkage being indirect, the inter-persona reward assumption, geographic coverage (20 of 100+ MultiTP countries; sparse personas for under-sampled countries), WVS Wave~7 vintage, and 4-bit/Int8 quantisation artefacts on IS stability.
    \item[] Guidelines:
    \begin{itemize}
        \item The answer \answerNA{} means that the paper has no limitation while the answer \answerNo{} means that the paper has limitations, but those are not discussed in the paper. 
        \item The authors are encouraged to create a separate ``Limitations'' section in their paper.
        \item The paper should point out any strong assumptions and how robust the results are to violations of these assumptions (e.g., independence assumptions, noiseless settings, model well-specification, asymptotic approximations only holding locally). The authors should reflect on how these assumptions might be violated in practice and what the implications would be.
        \item The authors should reflect on the scope of the claims made, e.g., if the approach was only tested on a few datasets or with a few runs. In general, empirical results often depend on implicit assumptions, which should be articulated.
        \item The authors should reflect on the factors that influence the performance of the approach. For example, a facial recognition algorithm may perform poorly when image resolution is low or images are taken in low lighting. Or a speech-to-text system might not be used reliably to provide closed captions for online lectures because it fails to handle technical jargon.
        \item The authors should discuss the computational efficiency of the proposed algorithms and how they scale with dataset size.
        \item If applicable, the authors should discuss possible limitations of their approach to address problems of privacy and fairness.
        \item While the authors might fear that complete honesty about limitations might be used by reviewers as grounds for rejection, a worse outcome might be that reviewers discover limitations that aren't acknowledged in the paper. The authors should use their best judgment and recognize that individual actions in favor of transparency play an important role in developing norms that preserve the integrity of the community. Reviewers will be specifically instructed to not penalize honesty concerning limitations.
    \end{itemize}

\item {\bf Theory assumptions and proofs}
    \item[] Question: For each theoretical result, does the paper provide the full set of assumptions and a complete (and correct) proof?
    \item[] Answer: \answerYes{}
    \item[] Justification: The single theoretical result, Theorem~\ref{thm:shrinkage} (minimax-optimality of disagreement-driven shrinkage in the James--Stein sense for panels of $N\geq 4$), is stated in the main text with all assumptions (panel size, Gaussian persona shifts, $\ell_2$ risk). The full proof, including referenced lemmas, is given in Appendix~\ref{app:bias_variance}. A short proof sketch and the connection to Eq.~\ref{eq:util_total} appear in \S\ref{sec:method}.
    \item[] Guidelines:
    \begin{itemize}
        \item The answer \answerNA{} means that the paper does not include theoretical results. 
        \item All the theorems, formulas, and proofs in the paper should be numbered and cross-referenced.
        \item All assumptions should be clearly stated or referenced in the statement of any theorems.
        \item The proofs can either appear in the main paper or the supplemental material, but if they appear in the supplemental material, the authors are encouraged to provide a short proof sketch to provide intuition. 
        \item Inversely, any informal proof provided in the core of the paper should be complemented by formal proofs provided in appendix or supplemental material.
        \item Theorems and Lemmas that the proof relies upon should be properly referenced. 
    \end{itemize}

    \item {\bf Experimental result reproducibility}
    \item[] Question: Does the paper fully disclose all the information needed to reproduce the main experimental results of the paper to the extent that it affects the main claims and/or conclusions of the paper (regardless of whether the code and data are provided or not)?
    \item[] Answer: \answerYes{}
    \item[] Justification: \S\ref{sec:method} gives the full algorithm (Algorithm~1, Eqs.~\ref{eq:util_total}--\ref{eq:disca_rel}); the full hyperparameter table (Table~\ref{tab:hyperparams}) lists every numerical default with the validation procedure (synthetic 200-scenario pool, transferability check on a 50-scenario MultiTP English subset). All backbones are publicly available open-weight checkpoints. Preprocessing (cap=80, oversample=36), seeds (3), and per-country preprocessing details are documented in Appendices~\ref{app:dataset_sens},~\ref{sec:multiseed}. We additionally release a self-contained reference implementation (\texttt{DISCA\_Alignment/}) with the same defaults and a single-command sweep entry point.
    \item[] Guidelines:
    \begin{itemize}
        \item The answer \answerNA{} means that the paper does not include experiments.
        \item If the paper includes experiments, a \answerNo{} answer to this question will not be perceived well by the reviewers: Making the paper reproducible is important, regardless of whether the code and data are provided or not.
        \item If the contribution is a dataset and\slash or model, the authors should describe the steps taken to make their results reproducible or verifiable. 
        \item Depending on the contribution, reproducibility can be accomplished in various ways. For example, if the contribution is a novel architecture, describing the architecture fully might suffice, or if the contribution is a specific model and empirical evaluation, it may be necessary to either make it possible for others to replicate the model with the same dataset, or provide access to the model. In general. releasing code and data is often one good way to accomplish this, but reproducibility can also be provided via detailed instructions for how to replicate the results, access to a hosted model (e.g., in the case of a large language model), releasing of a model checkpoint, or other means that are appropriate to the research performed.
        \item While NeurIPS does not require releasing code, the conference does require all submissions to provide some reasonable avenue for reproducibility, which may depend on the nature of the contribution. For example
        \begin{enumerate}
            \item If the contribution is primarily a new algorithm, the paper should make it clear how to reproduce that algorithm.
            \item If the contribution is primarily a new model architecture, the paper should describe the architecture clearly and fully.
            \item If the contribution is a new model (e.g., a large language model), then there should either be a way to access this model for reproducing the results or a way to reproduce the model (e.g., with an open-source dataset or instructions for how to construct the dataset).
            \item We recognize that reproducibility may be tricky in some cases, in which case authors are welcome to describe the particular way they provide for reproducibility. In the case of closed-source models, it may be that access to the model is limited in some way (e.g., to registered users), but it should be possible for other researchers to have some path to reproducing or verifying the results.
        \end{enumerate}
    \end{itemize}


\item {\bf Open access to data and code}
    \item[] Question: Does the paper provide open access to the data and code, with sufficient instructions to faithfully reproduce the main experimental results, as described in supplemental material?
    \item[] Answer: \answerYes{}
    \item[] Justification: We release \texttt{DISCA\_Alignment/} (anonymised for review; MIT-licensed) as supplemental material. The package is a self-contained slice of the research codebase: a single CLI entry point (\texttt{run.py}) drives the per-model 20-country sweep, with backend selection (\texttt{vllm} / \texttt{hf\_native} / \texttt{unsloth}) and explicit flags for the three external data paths (\texttt{--multitp-data}, \texttt{--wvs-data}, \texttt{--human-amce}). The README documents install (Python 3.10+, CUDA 12.x), HF token wiring for gated models, exact commands, and the output schema (per-country baseline / DISCA CSVs and a comparison file). MultiTP, WVS-7 Wave, and the human AMCE table are publicly available datasets cited in \S\ref{sec:exp:data}; access instructions are reproduced in the README.
    \item[] Guidelines:
    \begin{itemize}
        \item The answer \answerNA{} means that paper does not include experiments requiring code.
        \item Please see the NeurIPS code and data submission guidelines (\url{https://neurips.cc/public/guides/CodeSubmissionPolicy}) for more details.
        \item While we encourage the release of code and data, we understand that this might not be possible, so \answerNo{} is an acceptable answer. Papers cannot be rejected simply for not including code, unless this is central to the contribution (e.g., for a new open-source benchmark).
        \item The instructions should contain the exact command and environment needed to run to reproduce the results. See the NeurIPS code and data submission guidelines (\url{https://neurips.cc/public/guides/CodeSubmissionPolicy}) for more details.
        \item The authors should provide instructions on data access and preparation, including how to access the raw data, preprocessed data, intermediate data, and generated data, etc.
        \item The authors should provide scripts to reproduce all experimental results for the new proposed method and baselines. If only a subset of experiments are reproducible, they should state which ones are omitted from the script and why.
        \item At submission time, to preserve anonymity, the authors should release anonymized versions (if applicable).
        \item Providing as much information as possible in supplemental material (appended to the paper) is recommended, but including URLs to data and code is permitted.
    \end{itemize}


\item {\bf Experimental setting/details}
    \item[] Question: Does the paper specify all the training and test details (e.g., data splits, hyperparameters, how they were chosen, type of optimizer) necessary to understand the results?
    \item[] Answer: \answerYes{}
    \item[] Justification: There is no training (the method is inference-time on frozen weights); we therefore detail all decoding and controller settings. \S\ref{sec:experiments} specifies the dataset (MultiTP, 20 countries), preprocessing (deduplication, category cap=80, oversample to $\geq$36 per dimension), the seven main backbones with sizes, and the metric definitions (MIS, Pearson $r$, JSD, win rate). All hyperparameters and how they were chosen are tabulated in Table~\ref{tab:hyperparams} with the synthetic-pool grid-search procedure described in Appendix~\ref{app:hyperparams}. Sensitivity sweeps for the temperature family (Appendix~\ref{app:temp_sensitivity}), category cap (Appendix~\ref{app:r2_oversampling}), and the controller knobs ($s$, $\lambda_{\text{coop}}$, $\sigma$, $T_{\text{cat}}$) appear in Appendix~\ref{app:hyperparams}.
    \item[] Guidelines:
    \begin{itemize}
        \item The answer \answerNA{} means that the paper does not include experiments.
        \item The experimental setting should be presented in the core of the paper to a level of detail that is necessary to appreciate the results and make sense of them.
        \item The full details can be provided either with the code, in appendix, or as supplemental material.
    \end{itemize}

\item {\bf Experiment statistical significance}
    \item[] Question: Does the paper report error bars suitably and correctly defined or other appropriate information about the statistical significance of the experiments?
    \item[] Answer: \answerYes{}
    \item[] Justification: The headline table (Table~\ref{tab:main_macro_summary}) reports macro MIS as mean $\pm$ standard deviation across 3 independent seeds (independent IS noise draws; the underlying model weights and data are fixed). A separate multi-seed stability analysis (Appendix~\ref{sec:multiseed}) gives per-country variability, and a bootstrap confidence interval on JSD ($\pm 0.004$, Appendix~\ref{app:robustness_summary}) quantifies the noise floor for distributional metrics. We explicitly state that conclusions are stable well within the bootstrap noise floor; sensitivity sweeps along nine independent axes confirm this.
    \item[] Guidelines:
    \begin{itemize}
        \item The answer \answerNA{} means that the paper does not include experiments.
        \item The authors should answer \answerYes{} if the results are accompanied by error bars, confidence intervals, or statistical significance tests, at least for the experiments that support the main claims of the paper.
        \item The factors of variability that the error bars are capturing should be clearly stated (for example, train/test split, initialization, random drawing of some parameter, or overall run with given experimental conditions).
        \item The method for calculating the error bars should be explained (closed form formula, call to a library function, bootstrap, etc.)
        \item The assumptions made should be given (e.g., Normally distributed errors).
        \item It should be clear whether the error bar is the standard deviation or the standard error of the mean.
        \item It is OK to report 1-sigma error bars, but one should state it. The authors should preferably report a 2-sigma error bar than state that they have a 96\% CI, if the hypothesis of Normality of errors is not verified.
        \item For asymmetric distributions, the authors should be careful not to show in tables or figures symmetric error bars that would yield results that are out of range (e.g., negative error rates).
        \item If error bars are reported in tables or plots, the authors should explain in the text how they were calculated and reference the corresponding figures or tables in the text.
    \end{itemize}

\item {\bf Experiments compute resources}
    \item[] Question: For each experiment, does the paper provide sufficient information on the computer resources (type of compute workers, memory, time of execution) needed to reproduce the experiments?
    \item[] Answer: \answerYes{}
    \item[] Justification: Appendix~\ref{app:trigger} (Table~\ref{tab:r2_latency_compact}) reports per-scenario wall-clock on a single H100 GPU: 0.083--0.086 s/scenario for DISCA at $K\in\{64,128,256\}$ vs.\ 0.023 s/scenario for vanilla decoding (3.57--3.72$\times$ overhead). Backbone-level cost is given in Figure~\ref{fig:latency_vs_mis}: 350~ms/scenario for Phi-4 vs.\ 1414~ms for Llama-3.3-70B with DISCA. Sub-70B backbones run on a single H100 in BF16 via vLLM; the 70B backbone uses 4-bit Unsloth on the same hardware. The released code's CLI also runs end-to-end on a Kaggle T4/P100 dual-GPU node for sub-14B models. We acknowledge in Appendix~\ref{app:engineering} that exploratory experiments (25 controller variants on a 3-backbone $\times$ 5-country prototyping panel) consumed substantially more compute than the reported runs.
    \item[] Guidelines:
    \begin{itemize}
        \item The answer \answerNA{} means that the paper does not include experiments.
        \item The paper should indicate the type of compute workers CPU or GPU, internal cluster, or cloud provider, including relevant memory and storage.
        \item The paper should provide the amount of compute required for each of the individual experimental runs as well as estimate the total compute. 
        \item The paper should disclose whether the full research project required more compute than the experiments reported in the paper (e.g., preliminary or failed experiments that didn't make it into the paper). 
    \end{itemize}
    
\item {\bf Code of ethics}
    \item[] Question: Does the research conducted in the paper conform, in every respect, with the NeurIPS Code of Ethics \url{https://neurips.cc/public/EthicsGuidelines}?
    \item[] Answer: \answerYes{}
    \item[] Justification: The authors have reviewed the NeurIPS Code of Ethics. The work uses only public, ethically-collected datasets (MultiTP~\citep{jin2025multitp}, World Values Survey Wave~7, Moral Machine human AMCEs~\citep{awad2018moral}); it conducts no new human-subjects research and introduces no scraped data. Anonymity is preserved in the submission. Negative-impact considerations (risk of encoding harmful majorities) are explicitly discussed in \S\ref{sec:discussion} with a per-persona utility floor mitigation validated in Appendix~\ref{app:hyperparams} (Table~\ref{tab:r2_persona_floor}).
    \item[] Guidelines:
    \begin{itemize}
        \item The answer \answerNA{} means that the authors have not reviewed the NeurIPS Code of Ethics.
        \item If the authors answer \answerNo, they should explain the special circumstances that require a deviation from the Code of Ethics.
        \item The authors should make sure to preserve anonymity (e.g., if there is a special consideration due to laws or regulations in their jurisdiction).
    \end{itemize}


\item {\bf Broader impacts}
    \item[] Question: Does the paper discuss both potential positive societal impacts and negative societal impacts of the work performed?
    \item[] Answer: \answerYes{}
    \item[] Justification: \S\ref{sec:discussion} discusses both directions. Positive: a train-free, low-overhead controller lets practitioners respect cross-national preference heterogeneity at deployment time without finetuning per country, lowering the barrier to culturally pluralistic alignment. Negative: aligning to majority survey statistics could entrench harmful majorities or marginalise minority views~\citep{zewail2025moral}, and alignment to a survey statistic is not equivalent to legitimacy~\citep{atari2023humans}. We propose and validate a per-persona utility floor that caps how far any single persona's post-correction utility can drop from vanilla (Appendix~\ref{app:hyperparams}, Table~\ref{tab:r2_persona_floor}), and we are explicit that DISCA is a research artefact, not a deployment-ready normative system.
    \item[] Guidelines:
    \begin{itemize}
        \item The answer \answerNA{} means that there is no societal impact of the work performed.
        \item If the authors answer \answerNA{} or \answerNo, they should explain why their work has no societal impact or why the paper does not address societal impact.
        \item Examples of negative societal impacts include potential malicious or unintended uses (e.g., disinformation, generating fake profiles, surveillance), fairness considerations (e.g., deployment of technologies that could make decisions that unfairly impact specific groups), privacy considerations, and security considerations.
        \item The conference expects that many papers will be foundational research and not tied to particular applications, let alone deployments. However, if there is a direct path to any negative applications, the authors should point it out. For example, it is legitimate to point out that an improvement in the quality of generative models could be used to generate Deepfakes for disinformation. On the other hand, it is not needed to point out that a generic algorithm for optimizing neural networks could enable people to train models that generate Deepfakes faster.
        \item The authors should consider possible harms that could arise when the technology is being used as intended and functioning correctly, harms that could arise when the technology is being used as intended but gives incorrect results, and harms following from (intentional or unintentional) misuse of the technology.
        \item If there are negative societal impacts, the authors could also discuss possible mitigation strategies (e.g., gated release of models, providing defenses in addition to attacks, mechanisms for monitoring misuse, mechanisms to monitor how a system learns from feedback over time, improving the efficiency and accessibility of ML).
    \end{itemize}
    
\item {\bf Safeguards}
    \item[] Question: Does the paper describe safeguards that have been put in place for responsible release of data or models that have a high risk for misuse (e.g., pre-trained language models, image generators, or scraped datasets)?
    \item[] Answer: \answerNA{}
    \item[] Justification: We do not release any pretrained models, generative models, or scraped datasets. The released artefact (\texttt{DISCA\_Alignment/}) is an inference-time controller layered over publicly available open-weight LLMs whose existing safeguards remain in place. The method itself produces no novel generative capability beyond that of the underlying models; the per-persona utility floor (Appendix~\ref{app:hyperparams}, Table~\ref{tab:r2_persona_floor}) is included as an in-method safeguard against majoritarian skew.
    \item[] Guidelines:
    \begin{itemize}
        \item The answer \answerNA{} means that the paper poses no such risks.
        \item Released models that have a high risk for misuse or dual-use should be released with necessary safeguards to allow for controlled use of the model, for example by requiring that users adhere to usage guidelines or restrictions to access the model or implementing safety filters. 
        \item Datasets that have been scraped from the Internet could pose safety risks. The authors should describe how they avoided releasing unsafe images.
        \item We recognize that providing effective safeguards is challenging, and many papers do not require this, but we encourage authors to take this into account and make a best faith effort.
    \end{itemize}

\item {\bf Licenses for existing assets}
    \item[] Question: Are the creators or original owners of assets (e.g., code, data, models), used in the paper, properly credited and are the license and terms of use explicitly mentioned and properly respected?
    \item[] Answer: \answerYes{}
    \item[] Justification: All third-party assets are cited with version information. Datasets: MultiTP~\citep{jin2025multitp}, the Moral Machine human AMCE table~\citep{awad2018moral}, and World Values Survey Wave~7 (\texttt{WVS\_Cross-National\_Wave\_7\_inverted\_csv\_v6\_0.csv}) are used under the terms of the respective public releases. Models: Phi-4 and Phi-3.5-mini (MIT), Llama-3.3-70B (Llama-3 Community License), Qwen-2.5/3-VL (Apache-2.0/Qwen license), Gemma family (Gemma terms of use), Magistral-Small-2509 (Apache-2.0). Code dependencies (Transformers, vLLM, Unsloth, bitsandbytes) are listed with version pins in the released \texttt{requirements.txt}; we use them under their stated licenses (mostly Apache-2.0/MIT). The \texttt{compute\_ACME} reference implementation we follow is credited to~\citet{jin2025multitp} in Appendix~\ref{app:amce}.
    \item[] Guidelines:
    \begin{itemize}
        \item The answer \answerNA{} means that the paper does not use existing assets.
        \item The authors should cite the original paper that produced the code package or dataset.
        \item The authors should state which version of the asset is used and, if possible, include a URL.
        \item The name of the license (e.g., CC-BY 4.0) should be included for each asset.
        \item For scraped data from a particular source (e.g., website), the copyright and terms of service of that source should be provided.
        \item If assets are released, the license, copyright information, and terms of use in the package should be provided. For popular datasets, \url{paperswithcode.com/datasets} has curated licenses for some datasets. Their licensing guide can help determine the license of a dataset.
        \item For existing datasets that are re-packaged, both the original license and the license of the derived asset (if it has changed) should be provided.
        \item If this information is not available online, the authors are encouraged to reach out to the asset's creators.
    \end{itemize}

\item {\bf New assets}
    \item[] Question: Are new assets introduced in the paper well documented and is the documentation provided alongside the assets?
    \item[] Answer: \answerYes{}
    \item[] Justification: We release the DISCA reference implementation (\texttt{DISCA\_Alignment/}) as a new code asset under the MIT License. The package ships with: a top-level README documenting installation, backend selection (vLLM / HF native / Unsloth), the three external data flags, environment variables, the quick-start command, and the output schema; a flat \texttt{disca/} module with one-purpose files (controller, dual-pass core, runner, AMCE/metrics, persona generation, multilingual scenarios, model loader); pinned dependencies in \texttt{requirements.txt}; and a single CLI entry point (\texttt{run.py --help}). The submission archive is anonymised. All hyperparameter defaults match Table~\ref{tab:hyperparams} in the main paper.
    \item[] Guidelines:
    \begin{itemize}
        \item The answer \answerNA{} means that the paper does not release new assets.
        \item Researchers should communicate the details of the dataset\slash code\slash model as part of their submissions via structured templates. This includes details about training, license, limitations, etc. 
        \item The paper should discuss whether and how consent was obtained from people whose asset is used.
        \item At submission time, remember to anonymize your assets (if applicable). You can either create an anonymized URL or include an anonymized zip file.
    \end{itemize}

\item {\bf Crowdsourcing and research with human subjects}
    \item[] Question: For crowdsourcing experiments and research with human subjects, does the paper include the full text of instructions given to participants and screenshots, if applicable, as well as details about compensation (if any)?
    \item[] Answer: \answerNA{}
    \item[] Justification: The paper conducts no new crowdsourcing or human-subjects research. Human reference data is reused from the existing public Moral Machine experiment~\citep{awad2018moral} as packaged by MultiTP~\citep{jin2025multitp}, and from the World Values Survey Wave~7 public microdata; we make no new contact with participants and recruit no annotators.
    \item[] Guidelines:
    \begin{itemize}
        \item The answer \answerNA{} means that the paper does not involve crowdsourcing nor research with human subjects.
        \item Including this information in the supplemental material is fine, but if the main contribution of the paper involves human subjects, then as much detail as possible should be included in the main paper. 
        \item According to the NeurIPS Code of Ethics, workers involved in data collection, curation, or other labor should be paid at least the minimum wage in the country of the data collector. 
    \end{itemize}

\item {\bf Institutional review board (IRB) approvals or equivalent for research with human subjects}
    \item[] Question: Does the paper describe potential risks incurred by study participants, whether such risks were disclosed to the subjects, and whether Institutional Review Board (IRB) approvals (or an equivalent approval/review based on the requirements of your country or institution) were obtained?
    \item[] Answer: \answerNA{}
    \item[] Justification: No new human-subjects research is conducted. The work is a secondary analysis of pre-existing public datasets (Moral Machine human AMCEs, MultiTP scenarios, WVS-7 microdata), each of which carries its own ethics review at the original collection point; no further IRB approval applies.
    \item[] Guidelines:
    \begin{itemize}
        \item The answer \answerNA{} means that the paper does not involve crowdsourcing nor research with human subjects.
        \item Depending on the country in which research is conducted, IRB approval (or equivalent) may be required for any human subjects research. If you obtained IRB approval, you should clearly state this in the paper. 
        \item We recognize that the procedures for this may vary significantly between institutions and locations, and we expect authors to adhere to the NeurIPS Code of Ethics and the guidelines for their institution. 
        \item For initial submissions, do not include any information that would break anonymity (if applicable), such as the institution conducting the review.
    \end{itemize}

\item {\bf Declaration of LLM usage}
    \item[] Question: Does the paper describe the usage of LLMs if it is an important, original, or non-standard component of the core methods in this research? Note that if the LLM is used only for writing, editing, or formatting purposes and does \emph{not} impact the core methodology, scientific rigor, or originality of the research, declaration is not required.
    %this research?
    \item[] Answer: \answerYes{}
    \item[] Justification: LLMs are central to the work in two ways, both fully described. (i) Frozen open-weight LLMs (Phi-4, Llama-3.3-70B, Magistral-24B, Qwen3-VL-8B, Qwen2.5-7B, Phi-3.5-mini, Gemma-4-E2B) are the \emph{subject} that DISCA steers; their identities, sizes, and decoding settings are documented in \S\ref{sec:exp:models} and Appendix~\ref{app:model_landscape}. (ii) For the open-ended extension (\S\ref{sec:method:openended}), an LLM parses free-form responses into a (choice, confidence) pair to recover a pseudo-logit-gap; this auxiliary use is documented in that section. A separate synthetic 200-scenario tuning pool generated by GPT-4 is disclosed in Appendix~\ref{app:hyperparams} along with its limitations. LLMs were not used to generate the paper's claims, results, or proofs.
    \item[] Guidelines:
    \begin{itemize}
        \item The answer \answerNA{} means that the core method development in this research does not involve LLMs as any important, original, or non-standard components.
        \item Please refer to our LLM policy in the NeurIPS handbook for what should or should not be described.
    \end{itemize}

\end{enumerate}